\begin{table}[!htbp]
    \centering
    \caption{Pemetaan tujuan artefak terhadap aspek evaluasi.}
    \label{tab:tab6.1_pemetaan_tujuan_artefak_dan_evaluasi}
    \begin{tabular}{p{0.16\textwidth}p{0.29\textwidth}p{0.38\textwidth}}
        \hline
        \textbf{Artefak} & \textbf{Tujuan yang dievaluasi} & \textbf{Aspek evaluasi} \\
        \hline
        T1 & Modul prediksi kadar glukosa & Galat prediksi, Clarke Error Grid, horizon, finger-stick, ketidakpastian, deteksi hipoglikemia \\
        T2 & RAG terkondisi-prediksi & Hit@k, MRR, nDCG, natural vs divergen, lintas-pasien, \textit{oracle}, ablasi \\
        T3 & Integrasi CDSS \textit{doctor-mediated} & \textit{Generation}, sitasi, \textit{fallback}, waktu respons, memori, \textit{usability}, evaluasi ahli \\
        \hline
    \end{tabular}
\end{table}
